% META: {"id": "TNSI-ALGO-COURS-05", "chapitre": "TNSI-ALGORITHMIQUE", "type_objet": "cours", "capacites_codes": ["C5"], "sources_inspiration": [], "mode_creation": "ex_nihilo", "status": "approved", "capacites": ["TNSI-ALGORITHMIQUE-C5"]}
\section{Diviser pour régner}

% ============================================================
% STRATE 1 — Principe et tri fusion
% ============================================================

\definition[1]{
La méthode \textbf{diviser pour régner} résout un problème en trois étapes : \textbf{diviser} l'instance en sous-instances plus petites de même nature, \textbf{régner} en résolvant récursivement chaque sous-instance, puis \textbf{combiner} les solutions partielles pour obtenir la solution du problème initial.
}

\margeAppui{
Le \textbf{tri fusion} (\textit{merge sort}) est l'exemple emblématique : diviser la liste en deux moitiés, trier récursivement chaque moitié, puis fusionner les deux listes triées en une seule.
}

\textbf{Exemple.}
\begin{python}
def fusionner(gauche, droite):
    resultat = []
    i, j = 0, 0
    while i < len(gauche) and j < len(droite):
        if gauche[i] <= droite[j]:
            resultat.append(gauche[i])
            i += 1
        else:
            resultat.append(droite[j])
            j += 1
    return resultat + gauche[i:] + droite[j:]

def tri_fusion(liste):
    if len(liste) <= 1:
        return liste
    milieu = len(liste) // 2
    gauche = tri_fusion(liste[:milieu])
    droite = tri_fusion(liste[milieu:])
    return fusionner(gauche, droite)
\end{python}

% BEGIN-VERIFY
% def fusionner(gauche, droite):
%     resultat = []
%     i, j = 0, 0
%     while i < len(gauche) and j < len(droite):
%         if gauche[i] <= droite[j]:
%             resultat.append(gauche[i])
%             i += 1
%         else:
%             resultat.append(droite[j])
%             j += 1
%     return resultat + gauche[i:] + droite[j:]
%
% def tri_fusion(liste):
%     if len(liste) <= 1:
%         return liste
%     milieu = len(liste) // 2
%     gauche = tri_fusion(liste[:milieu])
%     droite = tri_fusion(liste[milieu:])
%     return fusionner(gauche, droite)
%
% assert tri_fusion([5, 3, 8, 1, 9, 2]) == [1, 2, 3, 5, 8, 9]
% assert tri_fusion([]) == []
% assert tri_fusion([1]) == [1]
% END-VERIFY

% ============================================================
% STRATE 2 — Cout
% ============================================================

\propriete[Coût du tri fusion]{
À chaque niveau de récursion, la liste est divisée en deux, ce qui donne $\log_2 n$ niveaux ; à chaque niveau, la fusion de toutes les sous-listes coûte $n$ comparaisons au total. Le coût total du tri fusion est donc de l'ordre de $n \log_2 n$, bien meilleur que les $n^2$ comparaisons d'un tri par sélection ou par insertion sur de grandes listes.
}

\erreurFrequente{
\textbf{Oublier le cas de base.} Sans le test \lstinline{if len(liste) <= 1: return liste}, la récursion ne s'arrête jamais : diviser une liste vide donnerait deux listes vides, et l'appel récursif se répéterait indéfiniment (ou lèverait une erreur de profondeur de récursion maximale atteinte).
}

\approfondissement{
La \textbf{rotation d'image} est un autre exemple classique de diviser pour régner : on découpe l'image en quadrants, on fait pivoter chaque quadrant récursivement, puis on réassemble le résultat en permutant les quadrants entre eux.
}
